%&latex
\documentclass[a4paper]{article}

\frenchspacing

\usepackage[cp1250]{inputenc}
\usepackage[czech]{babel}

\usepackage{a4wide}
\usepackage{amsmath, amsthm, amssymb, amsfonts}
\usepackage[mathcal]{eucal}
\usepackage{graphicx}
\usepackage{url}
\usepackage{color}
\usepackage{wrapfig}
\usepackage{capt-of}
\usepackage{float}
\usepackage{ifthen}
\usepackage{subfig}
\usepackage[normalem]{ulem}


% sirka a vyska textu nastavena jako papir, vsechny okraje vynulovany a pridano 20pt na kazdou stranu
% horizontalni rozmery
\setlength{\textwidth}{\paperwidth}
\addtolength{\textwidth}{-40pt}
\addtolength{\hoffset}{-1in}
\addtolength{\hoffset}{20pt}
\setlength{\oddsidemargin}{0in}
\setlength{\marginparsep}{0in}
% vertikalni rozmery
\setlength{\textheight}{\paperheight}
\addtolength{\textheight}{-60pt}
\addtolength{\voffset}{-1in}
\addtolength{\voffset}{20pt}
\setlength{\topmargin}{0in}
\setlength{\headheight}{0in}
\setlength{\headsep}{0in}


%Obrazek na miste
%pouziti
%%\obrazeknahore{adresa}{popisek}{label}
\long\def\obrazeknahore#1#2#3 {

\begin{figure}[t]
    \centering
    \includegraphics[width=0.8\textwidth]{#1}
    
    \caption{#2}
    \label{#3}
    
\end{figure}

}


%==========================================
%PEKELNA MAKRA NA ZAROVNANI OBRAZKU DOPRAVA

\makeatletter


%tohle je makro, ktere mi dovoluje obtekani i u kratkych environmentu
%ABSOLUTNE nechapu, jak to funguje, ale funguje to
%viz http://tex.stackexchange.com/questions/26078/ 
\def\odrovnej{\@@par
\ifnum\@@parshape=\z@ \let\WF@pspars\@empty \fi % reset `parshape'
\global\advance\c@WF@wrappedlines-\prevgraf \prevgraf\z@
\ifnum\c@WF@wrappedlines<\tw@ \WF@finale \fi}

\makeatother



%---
%makro, co da obrazek doprava a ostatni text ho obteka
%(bez toho predchazejiciho makra to ale poradne nebeha)
%pouziti:
%\obrazekvpravo{adresa}{popisek}{label}{procento sirky}
\long\def\obrazekvpravo#1#2#3#4{

\setlength\intextsep{-20pt}

    \begin{wrapfigure}{r}{#4\textwidth}
      \begin{center}
          \vspace{-10pt}
          
        \includegraphics[width=#4\textwidth]{#1}
        \vspace{-10pt}
        
      \end{center}
      
      \caption{#2}
      \label{#3}
      
      
    \end{wrapfigure}

\setlength\intextsep{0pt}

    
}




%---
%makro pro pripady, kdy wrapfigure neco mrsi
%je to docela pekelne
%je nutne mu dat jak text vpravo, tak text vlevo
%a nevim, jestli bude 100% fungovat, ale doufam, ze jo

%pouziti:
%\obrazekvpravominipage{adresa}{popisek}{label}{procento sirky}{1 - procento sirky}{text vlevo}
\long\def\obrazekvpravominipage#1#2#3#4#5#6{

\noindent\begin{minipage}{#5\linewidth}
\vspace{0pt}
#6
\end{minipage}
\hspace{0.5cm}
\noindent\begin{minipage}{#4\linewidth}
\vspace{0pt}
\centering
\includegraphics[width=0.9\textwidth]{#1}
\captionof{figure}{#2}
\label{#3}
\end{minipage}

}

%KONEC PEKELNYCH MAKER
%=====================

\def\lebka{\includegraphics[width=3mm]{../common/lebka}}

% makra pro poznamku u vyrokove a predikatove logiky
\def\vl{ -- ve v�rokov� logice}
\def\pl{ -- v predik�tov� logice}


%Vacsina prostredi je dvojjazicne. V pripade, ze znenie napr pozorovania je pisane po slovensky, malo by byt po slovensky aj oznacenie.

\def\ifis#1#2{\ifthenelse{\equal{#1}{0}}{}{#2}}

%\newenvironment{e}[3]{\pagebreak[2]\noindent\textbf{\ifis{#2}{$\bigstar$}#1}\ifis{#3}{\emph{~(#3)}}\par\noindent\leftskip 10pt }{\odrovnej\par\bigskip}
\newenvironment{e}[3]{\pagebreak[2]\noindent\textbf{\ifis{#2}{\lebka~}#1}\ifis{#3}{\emph{~(#3)}}\ifis{#2}{~\lebka}\par\noindent\leftskip 10pt }{\odrovnej\par\bigskip}


%obecne prostredie, ktore ma vyuzitie pri specialnych odstavcoch ako (uloha, algoritmus...) aby nevzniklo dalsich x prostredi
\newenvironment{obecne}[1]{\pagebreak[2]\noindent\textbf{#1}\par\noindent\leftskip 10pt}{\odrovnej\par\bigskip}

\newenvironment{report}{\pagebreak[2]\noindent\textbf{Report}\em\par\noindent\leftskip 10pt}{\par\bigskip}

%\newenvironment{reportN}[1]{\pagebreak[2]\noindent\textbf{Report~}\emph{(#1)}\emph\par\noindent\leftskip 10pt}{\odrovnej\par\bigskip}
\newenvironment{reportN}[1]{\pagebreak[2]\noindent\textbf{Report~}\emph{(#1)}\em\par\noindent\leftskip 10pt}{\odrovnej\par\bigskip}

\newenvironment{penumerate}{
\begin{enumerate}
  \setlength{\itemsep}{1pt}
  \setlength{\parskip}{0pt}
  \setlength{\parsep}{0pt}
  %\setlength{\topsep}{200pt}
  \setlength{\partopsep}{200pt}
}{\end{enumerate}}

\def\pismenka{\numberedlistdepth=2} %pouzit, ked clovek chce opismenkovany zoznam...

\newenvironment{pitemize}{
\begin{itemize}
  \setlength{\itemsep}{1pt}
  \setlength{\parskip}{0pt}
  \setlength{\parsep}{0pt}
}{\end{itemize}}

\definecolor{gris}{gray}{0.95}
\newcommand{\ramcek}[2]{\begin{center}\fcolorbox{white}{gris}{\parbox{#1}{#2}}\end{center}\par}


\title{\LARGE U�ebn� texty k st�tn� bakal��sk� zkou�ce \\ Matematika \\ Skal�rn� sou�in}

\begin{document}


\setcounter{section}{9}
\section{Skal�rn� sou�in}

\begin{e}{Po�adavky}{0}{0}
\begin{pitemize}
        \item Vlastnosti v re�ln�m i komplexn�m p��pad�
        \item Norma
        \item Cauchy-Schwarzova nerovnost
        \item Kolmost
        \item Ortogon�ln� dopln�k a jeho vlastnosti
\end{pitemize}
\end{e}

\subsection{Vlastnosti v re�ln�m i komplexn�m p��pad�}
\begin{e}{Definice}{0}{0}
Nech� $V$ je vektorov� prostor nad $\mathbb{C}$. Potom zobrazen� (funkce) z kart�zsk�ho sou�inu $V \times V \rightarrow \mathbb{C}$, kter� dvojici vektor� $x$ a $y$ p�irad� ��slo $\left<x,y\right>$ se naz�v� \emph{skal�rn� sou�in}, pokud spl�uje n�sleduj�c� axiomy (pro v�echny $x, x', y \in V$ a $\alpha, \beta \in \mathbb{C}$):
\begin{penumerate}
        \item $\left<x,x\right> \ge 0, \,\, \left<x,x\right> =0 \Leftrightarrow x = 0 \hfill \textit{(positivn� definitnost)}$
        \item $\left<\alpha x+\beta x', y\right> = \alpha \left<x,y\right> + \beta \left<x',y\right> \hfill \textit{(bilinearita)}$
                \begin{penumerate}
                        \item $\left<\alpha x,y\right> = \alpha \left<x,y\right>$
                        \item $\left<x+x',y\right> = \left<x,y\right> + \left<x',y\right>$
                \end{penumerate}
        \item $\left<x,y\right> = \overline{\left<y,x\right>} \hfill \textit{(symetrie - komplexn� sdru�en�)}$ \\
                \begin{e}{Pozn�mka}{0}{0}
                        Pro $V'$ nad $\mathbb{R}$ a vektory $\forall x,y \in V'$: $\left<x,y\right>=\left<y,x\right>$
                \end{e}
\end{penumerate}

Skal�rn� sou�in zna��me: $\left<x,y\right>$, $\left<x|y\right>$, $x.y$ \dots
\end{e}

\begin{e}{Pozorov�n�}{0}{0}
\begin{pitemize}
        \item $\left<x,x\right>=\overline{\left<x,x\right>}$, tedy je nutn� re�ln� ($\in\mathbb{R}$) i pro skal�rn� sou�iny nad $\mathbb{C}$
        \item $\left<x,\alpha y\right> = \overline{\left<\alpha y,x\right>} = \overline{\alpha}.\overline{\left<y,x\right>} = \overline{\alpha}.\left<x,y\right>$
        \item Skal�rn� sou�in m��e nab�vat z�porn�ch hodnot
\end{pitemize}
\end{e}

\begin{e}{Definice}{0}{0}
Ekvivalentn� definice: \emph{Skal�rn� sou�in} je pozitivn� definitn� (1) biline�rn� forma (2). V $\mathbb{R}$ nav�c symetrick� (3). V $\mathbb{C}$ nav�c forma, jej� matice je hermitovsk� (3).
\end{e}

\begin{e}{P��klady}{0}{0}
\begin{pitemize}
        \item \uv{Standardn�} skal�rn� sou�in pro $\mathbb{C}^n, \mathbb{R}^n$: $$\left<x,y\right> = \sum_{i=1}^{n} x_i \overline{y_i}$$
        \item Jin� sou�in v $\mathbb{R}^n$ definovan� pomoc� regul�rn� matice $A$ ��du $n$
                $$\left<x,y\right>=x^T A^T A y \;\;\;\;\;\; (\textit{pozorov�n�: } \left<x,x\right> = x^T A^T A x = \sum_{i=1}^n (Ax)_i^2)$$
        \item Skal�rn� sou�in ve vektorov�m prostoru $C[a,b]$ (integrovateln�ch funkc� na intervalu $[a,b]$):
                $$\left<f,g\right> = \int_a^b f(x) g(x) dx$$
\end{pitemize}
\end{e}

\subsection{Norma}

\begin{e}{Definice}{0}{Norma}
\emph{Norma} na vektorov�m prostoru V (nad $\mathbb{R}$ nebo nad $\mathbb{C}$) je zobrazen� $V \rightarrow \mathbb{R}$, kter� p�irad� vektoru $x \in V$ ��slo $\|x\|$ a spl�uje axiomy:
\begin{penumerate}
        \item $\forall x \in V: \|x\| \ge 0, \|x\|=0 \Leftrightarrow x = 0$
        \item $\forall x \in V, \forall \alpha \in \mathbb{C} (\mathbb{R}): \|\alpha x\| = |\alpha|.\|x\|$
        \item $\forall x,y \in V: \|x\| \ge 0, \|x+y\| \le \|x\| + \|y\| \hfill \textit{(troj�heln�kov� nerovnost)}$
\end{penumerate}
Norma $\|x\|$ m� v�znam \uv{d�lky} vektoru $x$.
\end{e}

\begin{e}{Definice}{0}{Normovan� vekt. prostor}
Vektorov� prostor s n�jakou normou naz�v�me \emph{normovan�}.
\end{e}

\begin{e}{P��klady}{0}{0}
\begin{pitemize}
\item Norma ur�en� skal�rn�m sou�inem
        $$\|x\| = \sqrt{\left<x,x\right>}$$
        \begin{e}{D�kaz}{0}{0}
                (1), (2) plyne z axiom� skal�rn�ho sou�inu, (3):
                \begin{align*}
                        \left<x,y\right>^2 \le \left<x,x\right>\left<y,y\right> & \Rightarrow & \left<x,y\right> & \le & \sqrt{\left<x,x\right> \left<y,y\right>} \\
                        & \Leftrightarrow & \left<x,x\right>+\left<y,y\right>+2\left<x,y\right> & \le & (\sqrt{\left<x,x\right>} + \sqrt{\left<y,y\right>})^2 \\
                        & \Leftrightarrow & \left<x+y,x+y\right> & \le & (\sqrt{\left<x,x\right>} + \sqrt{\left<y,y\right>})^2 \\
                        & \Leftrightarrow & \|x+y\| & \le & \|x\| + \|y\|
                \end{align*}
                Kde prvn� nerovnost je d�sledek Cauchy-Swarzovy nerovnosti...

                Ze standardn�ho skal�rn�ho sou�inu na $\mathbb{R}^n$ dostaneme euklidovskou normu (tj. \uv{d�lku} vektoru podle Pythagorovy v�ty) a euklidovskou vzd�lenost (vzd�lenost bod� $u$ a $v$ je $\|u-v\|$). Ka�d� vektorov� prostor se skal�rn�m sou�inem $\left<.,.\right>$ je normovan�m vektorov�m prostorem ($\|x\|=\sqrt{\left<x,x\right>}$), tedy i metrick�m prostorem ($d(x,y)=\|x-y\|$) a tedy i topologick�m prostorem.
        \end{e}
\item $L_1$ norma na $\mathbb{R}^n$:
        $$\|x\| = \sum_{i=1}^{n} |x_i|$$
\item $L_2$ norma na $\mathbb{C}^n$ - Euklidovsk� norma:
        $$\|x\| = \sqrt{\sum_{i=1}^{n} x_i \overline{x_i}}$$
\item $L_p$ norma na $\mathbb{R}^n$:
        $$\|x\| = \sqrt[p]{\sum_{i=1}^{n} |x_i|^p}$$
\item $L_{\infty}$ norma (stejn� jako $L_1$ norma neodpov�d� ��dn�mu skal�rn�mu sou�inu):
        $$\|x\| = \max_{i=1,\dots,n} \left( |x_i| \right)$$
\item Norma v prostoru integrovateln�ch funkc� na intervalu $[a,b]$ - $C[a,b]$
        $$\|f(x)\| = \int_a^b f^2(x) dx$$
\end{pitemize}
\end{e}


\subsection{Cauchy-Schwarzova nerovnost}


\begin{e}{V�ta}{0}{Cauchyho-Schwarzova nerovnost}
Nech� $V$ je prostor se skal�rn�m sou�inem nad $\mathbb{C}$ a $\|x\|$ je norma odvozen� ze skal�rn�ho sou�inu. Potom plat�:
$$\left|\left<x,y\right>\right| \le \|x\|\cdot\|y\| \;\;\;\;\;\; (\forall x,y \in V)$$

\begin{e}{D�kaz}{0}{0}
Pro $x=0$ nebo $y=0$ m�me $0�\le 0$.

Pro libovoln� $\alpha \in \mathbb{C}$ plat� $\|x+\alpha y\|^2 \ge 0$ (plat� i bez $( )^2$)
\begin{align*}
\|x+\alpha y\|^2 = \left<x+\alpha y, x+\alpha y\right> = \left<x, x+\alpha y\right> + \alpha \left<y, x+\alpha y\right> = \\
= \left<x,x\right> + \overline{\alpha}\left<x,y\right> + \alpha\left<y,x\right> + \alpha \overline{\alpha}\left<y,y\right>
\end{align*}
Zvol�me $\alpha = \frac{-\left<x,y\right>}{\left<y,y\right>}$ (\emph{t�m se eliminuj� $\overline{\alpha}\left<x,y\right>$ a $\alpha \overline{\alpha}\left<y,y\right>$})

Po dosazen�:
\begin{align*}
0 &\le& \left<x,x\right> + \alpha \left<y,x\right> \\
0 &\le& \left<x,x\right> - \frac{\left<x,y\right>}{\left<y,y\right>}\left<y,x\right> \\
\left<x,y\right>.\left<y,x\right> &\le& \left<x,x\right>.\left<y,y\right> \\
|\left<x,y\right>|^2 & \le & \|x\|^2 . \|y\|^2 \\
\textit{...a po odmocn�n�} \\
|\left<x,y\right>| & \le & \|x\|.\|y\|
\end{align*}

\end{e}

\begin{obecne}{Druh� mo�n� d�kaz}
\noindent Nadefinujeme prom�nnou $t\in\mathbb{R}$ a zavedeme funkci
$$p(t):=\left<u+t\cdot v,u+t\cdot v\right>=||u+tv||^2$$

\noindent
V�me: $p(t)\geq 0\ \forall t\in\mathbb{R}$ (z axiomu 1 skal. sou�inu). Z linearity plyne, �e $\left<u+tv,u+tv\right>=\left<u,u+tv\right>+t\left<v,u+tv\right>=\left<u,u\right>+t\left<u,v\right>+t\left<v,u\right>+t^2\left<v,v\right>=||u||^2+2t\left<u,v\right>+t^2||v||^2$. Tj. dost�v�me $p(t)$ jako kvadratickou funkci prom�nn� $t$:
$$p(t)=t^2||v||^2+2t\left<u,v\right>+||u||^2$$

\noindent
Proto�e $p(t)$ m� nez�porn� hodnoty na cel�m $\mathbb{R}$, mus� m�t tato rovnice max. jedno �e�en�, tj. diskriminant p�i po��t�n� ko�en� nesm� b�t kladn�:
$$D=b^2-4ac=4\left<u,v\right>^2-4||u||^2||v||^2\leq 0$$

\noindent
Po vyd�len� �ty�mi a odmocn�n� dost�v�me:
$$|\left<u,v\right>|\leq ||u||\cdot||v||$$
\end{obecne}

\begin{e}{D�sledek}{0}{0}
Platnost troj�heln�kov� nerovnosti pro normy odvozen� od skal�rn�ho sou�inu -- tj. normy odvozen� od skal�rn�ho sou�inu spl�uj� v�echny axiomy normy.
\end{e}

\begin{e}{D�sledek}{0}{0}
Nech� $x=(x_1,x_2,\dots,x_n)^T$, $y=(1,1,\dots,1)^T$ jsou dva vektory, pak pro standardn� skal�rn� sou�in plat�
\begin{align*}
|\left<x,y\right>| &= \sum_{i=1}^n x_i\cdot 1\\
\|x\| &= \sqrt{\sum_{i=1}^n x_i^2}\\
\|y\| &= \sqrt{n}
\end{align*}
po dosazen� do Cauchy-Schwarzovy nerovnosti okam�it� dostaneme nerovnost mezi aritmetick�m a kvadratick�m pr�m�rem
$$\frac{1}{n}\sum_{i=1}^{n} x_i \le \sqrt{\frac{1}{n}\sum_{i=1}^{n} x_i^2}$$
\end{e}



\begin{e}{D�sledek}{0}{0}
Ve vektorov�ch prostorech nad $\mathbb{R}$ a $\mathbb{C}$ lze definovat \emph{�hel}, sv�ran� dv�ma vektory: 
$$\cos\varphi = \frac{\left<u,v\right>}{\|u\|\ \|v\|}$$
a Cauchyho-Schwarzova nerovnost zaru�uje, �e $|\cos\varphi|\leq 1$.
\end{e}

\begin{e}{D�sledek}{0}{0}
Z takto definovan�ho �hlu mezi dv�ma vektory plyne i \emph{kosinov� v�ta}: 
$$\|u-v\|^2=\|u\|^2+\|v\|^2-2\|u\|\ \|v\|\cos\varphi$$
\end{e}
\end{e}


\subsection{Kolmost}
\begin{e}{Definice}{0}{kolm� vektory}
Vektory $x$ a�$y$�z prostoru se skal�rn�m sou�inem jsou vz�jemn� \emph{kolm�}�(\emph{ortogon�ln�}), pokud $\left<x,y\right>=0$, zna��me $x \bot y$.
\end{e}

\begin{e}{Definice}{0}{ortogon�ln� a ortonorm�ln� syst�m}
Soustava (syst�m) vektor� $v_1, \dots, v_n$ se naz�v� \emph{ortogon�ln�}, jestli�e $\left<v_i,v_j\right> = 0$ ($v_i \bot v_j$) pro $\forall i \neq j$ (tj. v�echny jej� vektory jsou navz�jem kolm�). 

\noindent
Plat�-li je�t� nav�c $\|v_i\|=1$ pro $\forall i=1,\dots,n$, jedn� se o soustavu \emph{ortonorm�ln�} (vektory jsou kolm� a nav�c maj� jednotkovou normu).
\end{e}

\begin{e}{Pozorov�n�}{0}{0}
Ka�d� syst�m nenulov�ch vz�jemn� kolm�ch vektor� (tj. i ortonorm�ln� nebo ortogon�ln�) je line�rn� nez�visl�.
\end{e}

\begin{e}{D�sledek}{0}{0}
Jestli�e ortogon�ln� syst�m generuje cel� vektorov� prostor, je jeho b�z�.
\end{e}

\begin{e}{Algoritmus}{0}{Gram-Schmidtova ortogonalizace}
Tento algoritmus zaji��uje p�eveden� libovoln� b�ze $(v_1,\dots,v_n)$ vektorov�ho prostoru $V$ na ekvivalentn� ortogon�ln� b�zi $(w_1,\dots,w_n)$. Ortonormalizace b�ze u� po jeho prob�hnut� znamen� jen vyn�soben� ka�d�ho $w_i$ ��slem $\frac{1}{\|w_i\|}$. Jeho pr�b�h:
\begin{penumerate}
    \item Zvolme $w_1 := v_1$.
    \item Pro $i$ postupn� od $1$ do $n$ opakujme:
\par\noindent
    Najdi $w_i=v_i-a_{i,1}w_1 - a_{i,2}w_2 - \dots - a_{i,i-1}w_{i-1}$ tak, aby pro $\forall j\in\{1,\dots,i\}$ platilo:
    $$w_i\bot w_j$$
    D� se uk�zat �e koeficienty $a_{i,j}$ jsou tvaru
    $$a_{i,j}=\frac{\left<v_i,w_j\right>}{\|w_j\|^2}$$
    \item Po $n$ iterac�ch dostaneme $w_1,\dots,w_n$ jako ortogon�ln� b�zi prostoru $V$.
\end{penumerate}

\textbf{Alternativn� postup - Gram-Schmidtova normalizace}:
\begin{penumerate}
        \item D�ny: $x_1,\dots,x_m \in V$ line�rn� nez�visl�.
        \item Pro $k=1,\dots,m$ prove�:
                $$y_k := x_k - \sum_{j=1}^{k-1} \left<x_k,z_j\right>z_j$$
                $$z_k := \frac{1}{\|y_k\|}y_k$$
        \item Ukon�i: $z_1,\dots,z_m$ je ortonorm�ln� syst�m ve $V$ a\\$\mathcal{L}(z_1,\dots,z_m) = \mathcal{L}(x_1,\dots,x_m)$
\end{penumerate}
\end{e}

\begin{e}{D�sledek}{0}{0}
Bu� $(v_1,\dots,v_n)$ b�ze vekt. prostoru se skal. sou�inem. Potom existuje ortonorm�ln� b�ze $(w_1,\dots,w_n)$, kdy pro ka�d� $k\in\{1,\dots,n\}$ je $\mathcal{L}(v_1,\dots,v_k)=\mathcal{L}(w_1,\dots,w_k)$. D�ky tomu se ka�d� ortogon�ln� syst�m vektor� v kone�n�dimension�ln�m vekt. prostoru se skal�rn�m sou�inem d� roz���it na ortogon�ln� b�zi (to m��eme d�ky Gram-Schmidtov� ortogonalizaci a Steinitzov� v�t� o v�m�n�).
\end{e}

\begin{e}{V�ta}{0}{Fourierovy koeficienty}
M�me-li danou n�jakou ortonorm�ln� b�zi $B=b_1,\dots,b_n$ vektorov�ho prostoru $V$, pak pro ka�d� $x\in V$ plat�:
$$x=\sum_{i=1}^n\left<x,b_i\right>b_i$$
a sou�adnice $\left<x,b_i\right>$ nazveme \emph{Fourierovy koeficienty} vektoru $x$.
\end{e}

\begin{e}{Pozn�mka}{0}{0}
Fourierovy �ady jsou sou�adnice funkc� ve vektorov�m prostoru spojit�ch funkc� na $[-\pi,\pi]$ se skal�rn�m sou�inem $\left<f,g\right>=\int_{-\pi}^{\pi}f(x)g(x)\mathrm{d}x$
\end{e}


\subsection{Ortogon�ln� dopln�k a jeho vlastnosti}
\begin{e}{Definice}{0}{0}
Nech� $V$ je mno�ina vektor� ve vektorov�m prostoru $W$ se skal�rn�m sou�inem. \emph{Ortogon�ln�m dopl�kem} $V$ (zna��me $V^{\bot}$) rozum�me mno�inu
$$V^{\bot}=\left\{v \in W; \forall x \in V: \left<v,x\right> = 0 \right\}$$
\end{e}

\begin{e}{Lemma}{0}{Vlastnosti}

Nech� $V$ je podprostor prostoru $W$ kone�n� dimenze. Potom plat�:
\begin{penumerate}
        \item $V^{\bot}$ je podprostor $W$
        \item $\dim(V^{\bot})=\dim(W)-\dim(V)$
        \item $(V^{\bot})^{\bot} = V \hfill \textit{(z roz�i�itelnosti ortogon�ln� b�ze)}$ 
        \item $V \cap V^{\bot} = \{0\}, \;\;\; V \oplus V^{\bot} = W \\ \textit{(operace} \oplus \textit{je spojen� dvou podprostor�...} \mathcal{L}(V \cup V^{\bot})\textit{)}$
        \item $U,V$ podprostory $W$. Je-li $U \subseteq V$, pak $U^{\bot} \supseteq V^{\bot}$\\
                $(x \in V^{\bot} \Leftrightarrow x \bot y \in V \Rightarrow x \bot u \in U \Leftrightarrow x \in U^{\bot})$
        \item $(U \cap V)^{\bot} = U^{\bot} \oplus V^{\bot}$
        \item $(U \oplus V)^{\bot} = U^{\bot} \cap V^{\bot}$
\end{penumerate}
\end{e}

\begin{e}{Definice}{0}{Ortogon�ln� projekce}
\emph{Ortogon�ln� projekce} vekt. prostoru $V$ na podprostor $U\subset V$ je zobrazen�, kter� ka�d�mu vektoru $v\in V$ p�i�ad� vektor $u\in U$ tak, �e 
$$\|v-u\|=\min\{\|v-w\|, w \in U\}$$
tedy vektor $u \in U$, kter� m� ze v�ech vektor� z $U$ nejmen�� vzd�lenost od $v$. Ten se pak naz�v� \emph{ortogon�ln� projekc�} vektoru $v$.
\end{e}
\newpage

\begin{e}{Report}{0}{Fiala} skal�rn� sou�in a kolmost, co� byl k�men �razu, proto�e o tom jsem toho moc nev�d�l, nicm�n�, "na to, abych v�s poslal dom� se to dou�it to nen�", tak�e za 3
\end{e}
\begin{e}{Report}{0}{Matou�ek} Skal�rn� sou�in: vlastnosti, norma, C-S nerovnost, a spol. (zadal Ml�ek, p�edv�d�l jsem Matou�kovi): Definice, jejich z�ejm� d�sledky v R a C, k �emu co slou��. Zeptal se jen co je ta ��rka nad (komplexn�m ��slem) a n�sledn� co je komplexn� sdru�en� ��slo, a cht�l rozepsat CS-nerovnost pro eukleidovsk� prostor. Zapomn�l jsem na jeden axiom skal�rn�ho sou�inu (nez�pornost $<u|u>$), to mi vytkl, a p�imhou�il oko i nad t�m,�e jsem mu zprvu tvrdil, �e plyne z ostatn�ch (spletl jsem si to s normou). Kolmost jsem stihl u� jen definovat, d�l jsem se nedostal. Je�t� jsem Matou�kovi tak trochu vnutil n�znak d�kazu CS nerovnosti (kdy� u� jsem se to u�il ;)), nebyl ale vy�adov�n. 
\end{e}
\begin{e}{Report}{0}{?} Skal�rn� sou�in, Cauchyho-Schwarzova nerovnost, Troj�heln�kov� nerovnost - zde se po m� vy�adoval d�kaz CS nerovnosti, co� jsem dal, + pak jsem m�l je�t� dok�zat tu troj�heln�kovou nerovnost z CS nerovnosti - tak� se mi to povedlo. Samoz�ejm� jsem tak� napsal axiomy skal�rn�ho sou�inu. za 1
\end{e}
\begin{e}{Report}{0}{Klazar} skal�rn� sou�in (SS, norma, C-S nerovnost, pouziti (ortogonalita)), co� jsem celkem i zvl�dnul a tak celkove cele statnice dopadly za 2:)
\end{e}
\begin{e}{Report}{0}{?} skal�rn� sou�in, norma, kolmost jsme toho moc ne�ek a co jsem �ekl, tak je�t� z p�lky �patn�. Hodnocen� matika za 3(s velkym m�nusem)
\end{e}
\begin{e}{Report}{0}{Majerech} Skal�rn� sou�in. Sta�ili v�etky defin�cie pojmov uveden�ch v zozname + p�r �al��ch vlastnost�. U CSB-nerovnosti som za�al robi� aj d�kaz, ale nepozeral na�.
\end{e}
\begin{e}{Report}{0}{Klazar} Skal�rn� prostory, Cauchyova-Schwarzova nerovnost (cht�l i jak souvis� s troj�heln�kovou nerovnost� normy), jinak definice a zn�n� sta�ily) 
\end{e}
\begin{e}{Report}{0}{Klazar} Skalarni soucin. Pak jsem zacala psat matiku, to uz bylo horsi, bo vsichni pomalu odchazeli, ja tam byla uz pekne dlouho a zaclo mne byt spatne a uz me to nebavilo. K skalarnimu soucinu jsem napsala definici, definici normy, Cauchy-Schwartzovu nerovnost + dukaz a pak par vec k ortogona a ortonormalizaci a Gramm-Schmitovu ortonormalizaci. To si jen precet a ze ok.
\end{e}
\begin{e}{Report}{0}{Fiala} Skalarni soucin, ortogonalita. Uplne vklidku, stacilo definovat pojmy, napsat neco kolem, pak polozil par praktictejsich otazek, doptal se na Fourriera vzhledem k ON bazi a proc se dela a hotovo. 
\end{e}
\begin{e}{Report}{0}{Fiala} skalarni soucin, ortogonalita. Ve druhym chtel k definici SS, ortogonality a fourierovych koeficientu slyset jeste neco jako, ze na $Z_n$ nema smysl SS zavadet, coz jsem mu nerekl a tak matika vysledna lepsi dvojka. 
\end{e}
\begin{e}{Report}{0}{Pangrac} Skalarny sucin & kolmost. Dokazat bolo treba len Cauchy Schwarzovu nerovnost, tu kolmost chcel len tak prehladovo. U dokazu tej CS nerovnosti som si nevedel spomenut na nejake technicke detaily, ale to myslim vobec nevadi. Celkovo sa nejde prilis do hlbky. Za jedna
\end{e}
\begin{e}{Report}{0}{�emli�ka} C-S nerovnost. Za�al jsem tou nerovnost�. Definoval jsem t�leso, VP, sk. sou�in a nakonec normu. Chyb�l mi jeden axiom a jin� vztah tam byl nav�c, proch�zen� a nal�z�n� bylo docela �morn�, d�kaz C-S nerovnosti jsem m�l z poloviny, ale kdy� prohl�sili, �e alfa je jen technick� z�le�itost, bylo to lep��, ne� cel� d�kaz �patn�. Cht�li d�sledky nerovnosti, t�eba jako troj�heln�kovou nebo srovn�n� pr�m�r�. Ten druh� jsme spole�n�mi silami odvodili, u� jsem byl op�t dost nerv�zn�.
\end{e}
\begin{e}{Report}{0}{Gregor} C-S nerovnost. Napsal jsem definice skalarniho soucinu a normy. U normy chtel nekolik moznych "typu" normy, tj. maximalni, manhattonska atd. coz jsem moc nevedel, musel jsem patrat v pameti jeste z dob kdy to bylo na prednasce. Pak jsem napsal CS nerovnost a ideu dukazu, na detaily jsem si nevzpomel. Potom trojuhelnikovou nerovnost. Tu jsem mel s dukazem pro R, ale on chtel i pro C. Nakonec jsem to s jeho pomoci dal do kupy. 
\end{e}
\begin{e}{Report}{0}{?} CS-nerovnost - bolo potreba definovat skalarny sucin a normu (dodatocne sa ma spytal, ci poznam nejaku normu, kt. nie je urcena skalarnym sucinom, tak som skusil tu s maximom) a potom som uviedol vetu CS-nerovnost + dokaz (spomenul som si na ten z vypracovanych otazok).
\end{e}
\begin{e}{Report}{0}{Fiala} Kolmost - �la l�pe, napsal jsem jim definice kolmosti, ortogon�ln�ho a ortonorm�ln�ho syst�mu a v�ty, �e ortogon�ln� syst�m vektor� je line�rn� nez�visl�, a tu o fourierov�ch sou�adnic�ch. K dopl�ku jsem krom� definice napsal v�tu o tom, �e j�dro matice je ortogon�ln� dopln�k jej�ho ��dkov�ho prostoru. V�e jsem p�edem dok�zal (je to jednoduch�, co� potom Fiala taky takto komentoval, ale proto�e jsem �ekal fiasko s �adami, tak jsem se rad�ji rozepsal v�c). Je�t� se ptali na n�jak� vlastnosti toho dopl�ku, Fiala cht�l hlavn� n�co o dimenz�ch, tak jsem mu napsal takov� to $dim(W dopln�k) + dim(W) = dim(V)$, je�t� ze m� vyt�hl jak z chlupat� deky, �e by se to dokazovalo n�jak p�es ortogon�ln� b�ze.
\end{e}


\end{document}
